
\documentclass[12pt,a4paper]{article}

\usepackage[a4paper,margin=1in]{geometry}
\usepackage{graphicx}
\usepackage{booktabs}
\usepackage{amsmath,amssymb}
\usepackage{float}
\usepackage{hyperref}
\usepackage{xcolor}
\usepackage{listings}
\usepackage{enumitem}
\usepackage{fancyhdr}
\usepackage{longtable}

\hypersetup{colorlinks=true,linkcolor=blue,urlcolor=blue}

\pagestyle{fancy}
\fancyhf{}
\lhead{ICS1512}
\rhead{Machine Learning Algorithms Laboratory}
\cfoot{\thepage}

\lstset{
language=Python,
basicstyle=\ttfamily\small,
numbers=left,
frame=single,
breaklines=true,
showstringspaces=false
}

\begin{document}

\begin{tabular}{|p{4cm}|p{11cm}|}
\hline
Experiment No. & 5 \\ \hline
Experiment Title & Decision Tree and Random Forest: A Comparative Classification Study
\footnote{\textbf{GitHub Repository: }\url{https://github.com/Barath-j593/ML-assignments}}\\ \hline
Student Name & NavinKumar S \\ \hline
Register Number & 3122247001039 \\ \hline
Faculty & Dr. Poreddy Ajay Kumar Reddy \\ \hline
Submission Date & 02/08/26 \\ \hline
\end{tabular}

\section{Objective}
To implement a Decision Tree classifier, extend it into a Random Forest ensemble model,
study the impact of hyperparameters on overfitting and generalization, select optimal
hyperparameters using 5-fold cross-validation, and compare the single-tree and ensemble-tree
models on the Wisconsin Diagnostic Breast Cancer dataset.

\section{Problem Statement}
Given the Wisconsin Diagnostic Breast Cancer dataset of 30 numerical features computed from
digitized images of breast mass fine-needle aspirates, the problem is binary classification:
predict whether a tumour is malignant or benign. The objective is to determine how tree depth
and other hyperparameters affect overfitting in a single Decision Tree, whether Random
Forest's bagging and random feature selection improve generalization over a single tree, and
which model is ultimately preferable for this diagnostic task.

\section{Dataset Description}
\begin{longtable}{|p{6cm}|p{8cm}|}
\hline
Dataset Name & Wisconsin Diagnostic Breast Cancer (WDBC) \\ \hline
Dataset Source & UCI Machine Learning Repository -- \url{https://archive.ics.uci.edu} (loaded via \texttt{sklearn.datasets.load\_breast\_cancer}) \\ \hline
Number of Samples & 569 \\ \hline
Number of Features & 30 numerical attributes (mean, standard error, and ``worst'' value of 10 cell-nucleus measurements) \\ \hline
Number of Classes & 2 (357 benign / 212 malignant) \\ \hline
Missing Values & 0 \\ \hline
Train-Test Split & 80:20 stratified (455 train / 114 test), \texttt{random\_state=42} \\ \hline
\end{longtable}

\section{Exploratory Data Analysis}
Include all relevant visualizations.
\begin{itemize}
\item Dataset overview
\item Statistical summary
\item Missing value analysis
\item Class distribution
\item Correlation matrix
\item Feature distribution
\end{itemize}

\subsection*{Class Distribution}
\begin{figure}[H]
\centering
\includegraphics[width=0.5\linewidth]{class_distribution.png}
\caption{Class distribution -- Breast Cancer Wisconsin (Diagnostic)}
\end{figure}
\textbf{Inference:} The dataset is mildly imbalanced (357 benign vs. 212 malignant, about
63:37). No missing values are present, so no missing-value visualization is required.

\subsection*{Correlation Matrix}
\begin{figure}[H]
\centering
\includegraphics[width=0.85\linewidth]{correlation_heatmap.png}
\caption{Correlation heatmap -- top 12 features correlated with diagnosis}
\end{figure}
\textbf{Inference:} ``Worst concave points'', ``worst perimeter'', ``worst radius'', and
``mean concave points'' show the strongest correlation with diagnosis (all $|r| > 0.7$),
consistent with the fact that larger, more irregularly-shaped nuclei are a hallmark of
malignancy. These same features later emerge as the top Random Forest feature importances
(Section 12) and the root splits of the visualized decision tree (Section 13), corroborating
each other.

\section{Data Preprocessing}
Explain:
\begin{itemize}
\item \textbf{Label encoding:} The target is already numerically encoded by scikit-learn
(0 = malignant, 1 = benign); no additional encoding step was required.
\item \textbf{Missing value handling:} Not required -- the dataset contains zero missing
values across all 569 rows and 30 features.
\item \textbf{Feature scaling:} Not required for either model -- Decision Trees and Random
Forests split on individual feature thresholds and are invariant to monotonic feature
scaling, unlike distance-based or gradient-based models.
\item \textbf{Train-test split:} 80:20 stratified split (455 train / 114 test rows),
preserving the 63:37 class ratio in both subsets.
\end{itemize}

\section{Machine Learning Algorithm}
Explain:
\begin{itemize}
\item \textbf{Working principle:} A Decision Tree recursively partitions the feature space by
selecting, at each node, the feature and threshold that most reduces class impurity (Gini
index or entropy), forming a hierarchy of if-then decision rules ending in leaf-node class
predictions. Random Forest is an ensemble of many such trees, each trained on a bootstrap
sample of the training data (bagging) and each considering only a random subset of features
at every split; final predictions are made by majority vote across all trees.
\item \textbf{Mathematical formulation:}
$$\text{Gini}(t) = 1 - \sum_{c} p(c\mid t)^2, \qquad
\text{Entropy}(t) = -\sum_{c} p(c\mid t)\log_2 p(c\mid t)$$
$$\hat{y} = \operatorname*{mode}_{i=1,\dots,B} T_i(x)$$
where $T_i$ is the $i$-th tree trained on a bootstrap sample, and $B$ is the number of trees
(\texttt{n\_estimators}).
\item \textbf{Advantages:} A single Decision Tree is highly interpretable (Figure~13) and
requires no feature scaling. Random Forest substantially reduces variance relative to a
single tree by averaging over many decorrelated trees (via bagging and random feature
selection), typically improving generalization.
\item \textbf{Limitations:} A single Decision Tree is prone to high variance and overfitting,
especially as depth grows unconstrained (Section 9). Random Forest loses the single tree's
direct interpretability (though feature importances partially compensate, Section 12),
requires more memory and training time to fit many trees, and its many hyperparameters
increase tuning cost (Section 8).
\end{itemize}

\section{Experimental Setup}
\begin{tabular}{ll}
\toprule
Python Version & \\
Scikit-Learn Version & \\
IDE & Jupyter Notebook \\
Operating System & \\
Hardware & \\
\bottomrule
\end{tabular}

\section{Hyperparameter Selection}

\begin{longtable}{|p{3.2cm}|p{2.6cm}|p{5.5cm}|p{2.3cm}|}
\hline
\textbf{Model} & \textbf{Search Method} & \textbf{Best Parameters} & \textbf{Best CV Accuracy} \\ \hline
\endhead
Decision Tree & Grid Search & criterion=gini, max\_depth=5, min\_samples\_split=2, min\_samples\_leaf=4 & 0.9385 \\ \hline
Random Forest & Grid Search & n\_estimators=100, max\_depth=10, max\_features=log2, bootstrap=True & 0.9648 \\ \hline
\end{longtable}

\textbf{Note:} Decision Tree was searched over criterion $\in \{gini, entropy\} \times
max\_depth \in \{2,3,4,5,7,10,None\} \times min\_samples\_split \in \{2,5,10\} \times
min\_samples\_leaf \in \{1,2,4\}$ (126 combinations). Random Forest was searched over
n\_estimators $\in \{50,100,200\} \times max\_depth \in \{3,5,10,None\} \times max\_features
\in \{sqrt, log2\} \times bootstrap \in \{True, False\}$ (48 combinations). Both used 5-fold
cross-validation optimizing accuracy.

\subsection*{Decision Tree Cross-Fold Results (summarized by criterion and max\_depth)}
\begin{longtable}{|p{3cm}|p{3cm}|p{3.5cm}|}
\hline
\textbf{Criterion} & \textbf{Max Depth} & \textbf{Avg CV Accuracy} \\ \hline
\endhead
gini & 2 & 0.9311 \\ \hline
gini & 3 & 0.9333 \\ \hline
gini & 5 & \textbf{0.9385} \\ \hline
gini & 10 & 0.9233 \\ \hline
gini & None & 0.9226 \\ \hline
entropy & 2 & 0.9209 \\ \hline
entropy & 3 & 0.9333 \\ \hline
entropy & 5 & 0.9263 \\ \hline
entropy & 10 & 0.9275 \\ \hline
entropy & None & 0.9250 \\ \hline
\end{longtable}
\textbf{Inference:} Accuracy is highest around \texttt{max\_depth=3--5} for both criteria and
noticeably lower once the tree is left unconstrained (\texttt{max\_depth=None}), directly
reflecting overfitting at high depth -- consistent with the depth-vs-accuracy curve in
Section~9. Gini and entropy perform similarly overall, with Gini marginally ahead at the
selected depth.

\subsection*{Random Forest Cross-Fold Results (summarized by n\_estimators, max\_depth, max\_features)}
\begin{longtable}{|p{2.3cm}|p{2.3cm}|p{2.7cm}|p{3cm}|}
\hline
\textbf{n\_estimators} & \textbf{Max Depth} & \textbf{Max Features} & \textbf{Avg CV Accuracy} \\ \hline
\endhead
50 & 10 & log2 & 0.9615 \\ \hline
50 & 10 & sqrt & 0.9582 \\ \hline
100 & 10 & log2 & \textbf{0.9626} \\ \hline
100 & 10 & sqrt & 0.9604 \\ \hline
100 & None & log2 & 0.9648 \\ \hline
200 & 10 & log2 & 0.9615 \\ \hline
200 & None & log2 & 0.9626 \\ \hline
\end{longtable}
\textbf{Inference:} Accuracy is fairly stable across n\_estimators once past 50 trees, and
\texttt{max\_features=log2} consistently edges out \texttt{sqrt} on this 30-feature dataset.
Unlike the Decision Tree, deeper/unconstrained Random Forest trees do not hurt generalization
-- bagging's variance reduction compensates for individual trees overfitting, which is the
central theoretical advantage of the ensemble.

\section{Results}

\begin{longtable}{|p{4.5cm}|p{2.3cm}|p{2.3cm}|}
\hline
\textbf{Metric} & \textbf{Decision Tree (tuned)} & \textbf{Random Forest (tuned)} \\ \hline
\endhead
Accuracy & 0.9035 & \textbf{0.9561} \\ \hline
Precision & 0.9420 & \textbf{0.9589} \\ \hline
Recall & 0.9028 & \textbf{0.9722} \\ \hline
F1-score & 0.9220 & \textbf{0.9655} \\ \hline
ROC-AUC & 0.9358 & \textbf{0.9924} \\ \hline
Training Time (s) & 0.0061 & 0.1489 \\ \hline
\end{longtable}

\section{Result Visualization}

Include all graphs generated during the experiment.

\begin{figure}[H]
\centering
\includegraphics[width=0.7\linewidth]{depth_overfitting_curve.png}
\caption{Decision Tree: effect of max\_depth on train/test accuracy}
\end{figure}
\textbf{Inference:} Training accuracy climbs to a perfect 1.0 by \texttt{max\_depth=7} and
stays there, while test accuracy peaks around \texttt{max\_depth=3--4} ($\approx$0.939) and
then \emph{declines} to a lower plateau ($\approx$0.912) as depth increases further -- a
textbook overfitting curve: the tree memorizes training noise at high depth without any
corresponding gain (in fact, a loss) in generalization.

\begin{figure}[H]
\centering
\includegraphics[width=0.6\linewidth]{n_estimators_curve.png}
\caption{Random Forest: effect of n\_estimators on test accuracy}
\end{figure}
\textbf{Inference:} Test accuracy rises from 0.947 at a single tree to a stable plateau of
0.956 from around 25 trees onward, with no overfitting as n\_estimators grows further --
unlike Decision Tree depth, adding more trees to a Random Forest does not increase variance,
since each additional tree only further averages out (rather than compounds) individual tree
error.

\begin{figure}[H]
\centering
\includegraphics[width=0.75\linewidth]{confusion_matrices.png}
\caption{Confusion matrices -- best Decision Tree vs. best Random Forest}
\end{figure}
\textbf{Inference:} Random Forest makes visibly fewer errors of both types (false positives
and false negatives) than the Decision Tree on the same test set, most notably fewer missed
malignant cases -- the clinically more costly error type.

\begin{figure}[H]
\centering
\includegraphics[width=0.55\linewidth]{roc_curves.png}
\caption{ROC curves -- best Decision Tree vs. best Random Forest}
\end{figure}
\textbf{Inference:} Random Forest's ROC curve dominates the Decision Tree's across almost the
entire range (AUC 0.992 vs. 0.936), confirming its probability-ranking quality is
substantially better, not just its hard classification accuracy.

\begin{figure}[H]
\centering
\includegraphics[width=0.55\linewidth]{cv_accuracy.png}
\caption{5-fold cross-validation accuracy -- tuned Decision Tree vs. tuned Random Forest}
\end{figure}
\textbf{Inference:} Random Forest outperforms Decision Tree on every one of the 5 folds
(average 0.9596 vs. 0.9245), with visibly less fold-to-fold variance -- direct empirical
confirmation of bagging's variance-reduction effect, not just a single-split artifact.

\begin{figure}[H]
\centering
\includegraphics[width=0.7\linewidth]{feature_importance.png}
\caption{Random Forest -- top 10 feature importances}
\end{figure}
\textbf{Inference:} ``Worst concave points'', ``worst area'', and ``worst radius'' are the
three most important features by a clear margin, matching the strongest correlations
identified in the EDA correlation heatmap (Section 4) and the root/near-root splits of the
visualized decision tree below -- three independent views of the data agreeing on the same
handful of dominant features.

\begin{figure}[H]
\centering
\includegraphics[width=0.9\linewidth]{tree_visualization.png}
\caption{Decision tree visualization (max\_depth=3, for interpretability)}
\end{figure}
\textbf{Inference:} The tree's root split is on ``worst radius'' ($\le 16.795$), immediately
separating a large, highly pure benign-leaning branch (266 of 271 samples benign) from a
malignant-leaning branch -- illustrating in a single readable diagram exactly why tree-based
models are valued for interpretability, a property the Random Forest ensemble trades away for
accuracy.

\begin{figure}[H]
\centering
\includegraphics[width=0.7\linewidth]{classifier_comparison.png}
\caption{Best Decision Tree vs. best Random Forest -- Accuracy / Precision / Recall / F1}
\end{figure}
\textbf{Inference:} Random Forest leads on every metric, with its largest margin over Decision
Tree on Recall (0.972 vs. 0.903) -- for a cancer-diagnosis task, this recall gap (fewer missed
malignant cases) is arguably the most clinically important difference between the two models.

\section{5-Fold Cross-Validation Performance Comparison}
\begin{longtable}{|p{2.3cm}|p{1.6cm}|p{1.6cm}|p{1.6cm}|p{1.6cm}|p{1.6cm}|p{1.8cm}|}
\hline
\textbf{Model} & \textbf{Fold 1} & \textbf{Fold 2} & \textbf{Fold 3} & \textbf{Fold 4} & \textbf{Fold 5} & \textbf{Average} \\ \hline
\endhead
Decision Tree & 0.9211 & 0.8947 & 0.9386 & 0.9298 & 0.9381 & 0.9245 \\ \hline
Random Forest & 0.9649 & 0.9561 & 0.9561 & 0.9561 & 0.9646 & \textbf{0.9596} \\ \hline
\end{longtable}

\section{Observation Questions}
\begin{enumerate}
\item \textbf{How does tree depth affect overfitting in Decision Trees?} As shown in
Figure~7, training accuracy reaches 100\% by depth 7 while test accuracy peaks at depth 3--4
and then declines -- deeper trees fit training-set-specific noise (individual sample
peculiarities) rather than generalizable patterns, the defining signature of overfitting.
\item \textbf{Which hyperparameter had the greatest impact on performance?} For the Decision
Tree, \texttt{max\_depth} had by far the largest effect (CV accuracy swinging from 0.921 to
0.939 across depths, with a clear optimum), more than criterion choice or leaf/split minimums.
For Random Forest, \texttt{max\_features} (log2 vs. sqrt) and using an unconstrained
\texttt{max\_depth} together produced the best result, though the overall spread across the
Random Forest grid (0.950--0.965) was much narrower than the Decision Tree's, indicating
Random Forest is considerably less hyperparameter-sensitive.
\item \textbf{How does Random Forest improve generalization?} By training many trees on
bootstrap-resampled data and restricting each split to a random feature subset, individual
trees become decorrelated from one another; averaging their predictions cancels out much of
each tree's individual overfitting/variance, which is exactly why Random Forest's
train--test gap does not widen with more trees (Figure~8) the way the Decision Tree's does
with depth (Figure~7).
\item \textbf{Did ensemble learning always improve performance? Why or why not?} Yes, in this
experiment Random Forest outperformed the tuned Decision Tree on every metric and every
cross-validation fold, with no exceptions. This is not a universal guarantee for all
datasets, however -- ensembling primarily helps when the base learner's error is dominated by
variance (as a Decision Tree's typically is); on a dataset where a single tree already
generalizes near-optimally with low variance, the ensemble's benefit would be much smaller.
\end{enumerate}

\section{Discussion}
The Wisconsin Diagnostic Breast Cancer dataset demonstrates textbook Decision Tree
overfitting behaviour: unconstrained trees reach perfect training accuracy while test
accuracy actually degrades past a shallow optimal depth. Random Forest's bagging and random
feature selection directly address this, delivering both higher accuracy (0.956 vs. 0.904)
and lower cross-validation variance (Section 14) at a modest training-time cost (0.149s vs.
0.006s -- still negligible in absolute terms at this dataset size). The strong agreement
between the correlation heatmap, Random Forest feature importances, and the visualized
decision tree's early splits (all highlighting ``worst radius,'' ``worst concave points,'' and
related ``worst''-prefixed measurements) gives confidence that both models are learning
genuine, clinically-sensible signal rather than spurious correlations. A limitation of this
experiment is that hyperparameter search used accuracy as the sole optimization target; for a
diagnostic task, optimizing directly for recall (minimizing missed malignant cases) may be
more clinically appropriate and could shift the selected hyperparameters.

\section{Conclusion}
For the Wisconsin Diagnostic Breast Cancer dataset, Random Forest (100 trees, max\_depth
unconstrained, max\_features=log2) is unambiguously the better-performing model, exceeding
the tuned Decision Tree on every evaluation metric and every cross-validation fold, most
notably on recall -- the most clinically consequential metric for a cancer-diagnosis task.
The Decision Tree's overfitting at high depth was clearly demonstrated and corrected via
hyperparameter tuning (optimal depth $\approx$5), but even the best-tuned single tree could
not match the ensemble's generalization. The trade-off is a loss of the single tree's direct
visual interpretability, partially offset by Random Forest's feature-importance ranking,
which converges on the same dominant features the tree itself split on first.

\section{References}
\begin{enumerate}[label={[\arabic*]}]
\item Scikit-learn Documentation, ``Decision Trees,'' \url{https://scikit-learn.org/stable/modules/tree.html}
\item Scikit-learn Documentation, ``Random Forests,'' \url{https://scikit-learn.org/stable/modules/ensemble.html#forest}
\item UCI Machine Learning Repository, ``Breast Cancer Wisconsin (Diagnostic) Data Set,'' \url{https://archive.ics.uci.edu}
\end{enumerate}

\end{document}
